%-----------------------------------------------------------------------------------------------------------------------------------------------%
%	The MIT License (MIT)
%
%	Copyright (c) 2015 Jan Küster
%
%	Permission is hereby granted, free of charge, to any person obtaining a copy
%	of this software and associated documentation files (the "Software"), to deal
%	in the Software without restriction, including without limitation the rights
%	to use, copy, modify, merge, publish, distribute, sublicense, and/or sell
%	copies of the Software, and to permit persons to whom the Software is
%	furnished to do so, subject to the following conditions:
%
%	THE SOFTWARE IS PROVIDED "AS IS", WITHOUT WARRANTY OF ANY KIND, EXPRESS OR
%	IMPLIED, INCLUDING BUT NOT LIMITED TO THE WARRANTIES OF MERCHANTABILITY,
%	FITNESS FOR A PARTICULAR PURPOSE AND NONINFRINGEMENT. IN NO EVENT SHALL THE
%	AUTHORS OR COPYRIGHT HOLDERS BE LIABLE FOR ANY CLAIM, DAMAGES OR OTHER
%	LIABILITY, WHETHER IN AN ACTION OF CONTRACT, TORT OR OTHERWISE, ARISING FROM,
%	OUT OF OR IN CONNECTION WITH THE SOFTWARE OR THE USE OR OTHER DEALINGS IN
%	THE SOFTWARE.
%
%
%-----------------------------------------------------------------------------------------------------------------------------------------------%


%============================================================================%
%
%	DOCUMENT DEFINITION
%
%============================================================================%

%we use article class because we want to fully customize the page and dont use a cv template
\documentclass[10pt]{article}


%----------------------------------------------------------------------------------------
%	ENCODING
%----------------------------------------------------------------------------------------

%we use utf8 since we want to build from any machine
\usepackage[utf8]{inputenc}

%----------------------------------------------------------------------------------------
%	LOGIC
%----------------------------------------------------------------------------------------

% provides \isempty test
\usepackage{xifthen}

%----------------------------------------------------------------------------------------
%	FONT
%----------------------------------------------------------------------------------------

% some tex-live fonts - choose your own

%\usepackage[defaultsans]{droidsans}
%\usepackage[default]{comfortaa}
%\usepackage{cmbright}
\usepackage[default]{raleway}
%\usepackage{fetamont}
%\usepackage[default]{gillius}
%\usepackage[light,math]{iwona}
%\usepackage[thin]{roboto}

% set font default
\renewcommand*\familydefault{\sfdefault}
\usepackage[T1]{fontenc}

% more font size definitions
\usepackage{moresize}


%----------------------------------------------------------------------------------------
%	PAGE LAYOUT  DEFINITIONS
%----------------------------------------------------------------------------------------

%debug page outer frames
%\usepackage{showframe}


%define page styles using geometry
\usepackage[letterpaper]{geometry}

% Adjust the margins - roughly 1/2 in all around
\geometry{top=0.5in, bottom=-0.24in, left=0.5in, right=0.5in}

%use customized header
\usepackage{fancyhdr}
\pagestyle{fancy}

%less space between header and content
\setlength{\headheight}{-5pt}


%customize entries left, center and right
%\lhead{}
%\chead{}
%\rhead{}


%indentation is zero
\setlength{\parindent}{0mm}

%----------------------------------------------------------------------------------------
%	TABLE /ARRAY DEFINITIONS
%----------------------------------------------------------------------------------------

%for layouting tables
\usepackage{multicol}
\usepackage{multirow}

%extended aligning of tabular cells
\usepackage{array}

\newcolumntype{x}[1]{%
>{\raggedleft\hspace{0pt}}p{#1}}%


%----------------------------------------------------------------------------------------
%	GRAPHICS DEFINITIONS
%----------------------------------------------------------------------------------------

%for header image
\usepackage{graphicx}

%for floating figures
\usepackage{wrapfig}
\usepackage{float}
%\floatstyle{boxed}
%\restylefloat{figure}

%for drawing graphics
\usepackage{tikz}
\usetikzlibrary{shapes, backgrounds,mindmap, trees}


%----------------------------------------------------------------------------------------
%	Color DEFINITIONS
%----------------------------------------------------------------------------------------

\usepackage{color}

%accent color
%\definecolor{sectcol}{RGB}{0,150,255}
\definecolor{sectcol}{RGB}{0,28,224}

%dark background color
%\definecolor{bgcol}{RGB}{110,110,110}
\definecolor{bgcol}{RGB}{64,64,64}

%light background / accent color
%\definecolor{softcol}{RGB}{225,225,225}
\definecolor{softcol}{RGB}{128,128,128}


%============================================================================%
%
%
%	DEFINITIONS
%
%
%============================================================================%

%----------------------------------------------------------------------------------------
% 	HEADER
%----------------------------------------------------------------------------------------

% remove top header line
\renewcommand{\headrulewidth}{0pt}

%remove botttom header line
\renewcommand{\footrulewidth}{0pt}

%remove pagenum
\renewcommand{\thepage}{}

%remove section num
\renewcommand{\thesection}{}

%----------------------------------------------------------------------------------------
% 	ARROW GRAPHICS in Tikz
%----------------------------------------------------------------------------------------

% a six pointed arrow poiting to the left
\newcommand{\tzlarrow}{(0,0) -- (0.2,0) -- (0.3,0.2) -- (0.2,0.4) -- (0,0.4) -- (0.1,0.2) -- cycle;}

% include the left arrow into a tikz picture
% param1: fill color
%
\newcommand{\larrow}[1]
{\begin{tikzpicture}[scale=0.58]
	 \filldraw[fill=#1!100,draw=#1!100!black]  \tzlarrow
 \end{tikzpicture}
}

% a six pointed arrow poiting to the right
\newcommand{\tzrarrow}{ (0,0.2) -- (0.1,0) -- (0.3,0) -- (0.2,0.2) -- (0.3,0.4) -- (0.1,0.4) -- cycle;}

% include the right arrow into a tikz picture
% param1: fill color
%
\newcommand{\rarrow}
{\begin{tikzpicture}[scale=0.7]
	\filldraw[fill=sectcol!100,draw=sectcol!100!black] \tzrarrow
 \end{tikzpicture}
}



%----------------------------------------------------------------------------------------
%	custom sections
%----------------------------------------------------------------------------------------

% create a coloured box with arrow and title as cv section headline
% param 1: section title
%
\newcommand{\cvsection}[1]
{
	\begin{center}
		\large\textcolor{sectcol}{\textbf{#1}}
	\end{center}
}

%create a coloured arrow with title as cv meta section section
% param 1: meta section title
%
\newcommand{\metasection}[2]
{
%\begin{tabular*}{1\textwidth}{r r}
\footnotesize{#2} \hspace*{\fill} \footnotesize{#1}\\[1pt]
%\end{tabular*}
}

%----------------------------------------------------------------------------------------
%	 CV EVENT
%----------------------------------------------------------------------------------------

% creates a stretched box as cv entry headline followed by two paragraphs about
% the work you did
% param 1:	event time i.e. 2014 or 2011-2014 etc.
% param 2:	event name (what did you do?)
% param 3:	institution (where did you work / study)
% param 4:	what was your position
%
\newcommand{\cveventone}[4]
{

\begin{tabular*}{1\textwidth}{p{6.15in}  x{1in}}
	\textbf{#2} - \textcolor{bgcol}{#3} &   \vspace{2.5pt}\textcolor{sectcol}{#1}
\end{tabular*}

\vspace{-8pt}
\textcolor{softcol}{\hrule}
\vspace{6pt}

  $\cdot$ #4\\[3pt]

}

% creates a stretched box as cv entry headline followed by two paragraphs about
% the work you did
% param 1:	event time i.e. 2014 or 2011-2014 etc.
% param 2:	event name (what did you do?)
% param 3:	institution (where did you work / study)
% param 4:	what was your position
% param 5:	some words about your contributions
%
%\newcommand{\cveventtwo}[5]
%{
%
%\begin{tabular*}{1\textwidth}{p{13.6cm}  x{3.9cm}}
%	\textbf{#2} - \textcolor{bgcol}{#3} &   \vspace{2.5pt}\textcolor{sectcol}{#1}
%\end{tabular*}
%
%\vspace{-8pt}
%\textcolor{softcol}{\hrule}
%\vspace{6pt}
%
%  $\cdot$ #4\\[3pt]
%  $\cdot$ #5\\[6pt]
%
%}

% creates a stretched box as cv entry headline followed by two paragraphs about
% the work you did
% param 1:	event time i.e. 2014 or 2011-2014 etc.
% param 2:	event name (what did you do?)
% param 3:	institution (where did you work / study)
% param 4:	what was your position
% param 5:	some words about your contributions
%
\newcommand{\cveventtwo}[5]
{

\begin{tabular*}{1\textwidth}{p{6.15in}  x{1in}}
	\textbf{#2} - \textcolor{bgcol}{#3} &   \vspace{2.5pt}\textcolor{sectcol}{#1}
\end{tabular*}

\vspace{-8pt}
\textcolor{softcol}{\hrule}
\vspace{6pt}

  $\cdot$ #4\\[3pt]
  $\cdot$ #5\\[6pt]

}

% creates a stretched box as cv entry headline followed by two paragraphs about
% the work you did
% param 1:	event time i.e. 2014 or 2011-2014 etc.
% param 2:	event name (what did you do?)
% param 3:	institution (where did you work / study)
% param 4:	what was your position
% param 5:	some words about your contributions
% param 6:	some words about your contributions
%
\newcommand{\cveventthree}[6]
{

\begin{tabular*}{1\textwidth}{p{6.15in}  x{1in}}
	\textbf{#2} - \textcolor{bgcol}{#3} &   \vspace{2.5pt}\textcolor{sectcol}{#1}
\end{tabular*}

\vspace{-8pt}
\textcolor{softcol}{\hrule}
\vspace{6pt}

  $\cdot$ #4\\[3pt]
  $\cdot$ #5\\[3pt]
  $\cdot$ #6\\[6pt]

}

% creates a stretched box as
\newcommand{\cveventmeta}[2]
{
	\mbox{\mystrut \hspace{87pt}\textit{#1}}\\
	#2
}

%----------------------------------------------------------------------------------------
% CUSTOM STRUT FOR EMPTY BOXES
%----------------------------------------- -----------------------------------------------
\newcommand{\mystrut}{\rule[-.3\baselineskip]{0pt}{\baselineskip}}

%----------------------------------------------------------------------------------------
% CUSTOM LOREM IPSUM
%----------------------------------------------------------------------------------------
\newcommand{\lorem}
{Lorem ipsum dolor sit amet, consectetur adipiscing elit. Donec a diam lectus.}



%============================================================================%
%
%
%
%	DOCUMENT CONTENT
%
%
%
%============================================================================%
\begin{document}


%use our custom fancy header definitions
\pagestyle{fancy}





%----------------------------------------------------------------------------------------
%	HEADER IMAGE
%----------------------------------------------------------------------------------------

%\begin{figure}[H]
%\begin{flushright}
%	\includegraphics[trim= 320 130 460 210,clip,width=0.2\linewidth]{myfoto.jpg}	%trimming relative to image size!
%\end{flushright}
%\end{figure}


%---------------------------------------------------------------------------------------
%	TITLE HEADLINE
%----------------------------------------------------------------------------------------
\vspace{-8pt}
\begin{center}
	\HUGE \textsc{Eric Gustafson}
          %\textcolor{sectcol}{\rule[-1mm]{1mm}{0.9cm}} \textsc{Resume}
          \\[4pt]
	\small Software Architect | Technology Leader | Cloud Native | DevOps | Monitoring | Networking
\end{center}



\vspace{6pt}


%---------------------------------------------------------------------------------------
%	META SECTION
%----------------------------------------------------------------------------------------
\metasection{Colorado Springs, CO, USA}
            {\textbf{Role:} Software Architect and Technical Leader}
\metasection{www.elfwerks.org/ericg}
            {\textbf{Fields:} Software Architecture and Development,
              Cloud, Monitoring, Networking, DevOps}
\metasection{ericg@elfwerks.org}
            {\textbf{Technologies:} Go, Python, Java, C/C++, Ansible,
              Chef, Nagios/Icinga, ELK, Hadoop, Kafka, RabbitMQ, ZeroMQ}
\metasection{+1.719.359.4112}
            {\textbf{Cloud:} Microservices, Kubernetes, Rancher,
              CoreOS, OpenStack, AWS, Google}
\vspace{-2pt}
\textcolor{softcol}{\hrule}
\vspace{6pt}

\normalsize

%---------------------------------------------------------------------------------------
%	SUMMARAY (optional)
%----------------------------------------------------------------------------------------
\vspace{-6pt}
\cvsection{Summary}
Agile, detail oriented, hands-on technology leader with excellent
communication skills.  Extensive experience in Cloud native
development and operation, DevOps, System Management, Monitoring,
Networking and Messaging.

%============================================================================%
%
%	CV SECTIONS AND EVENTS (MAIN CONTENT)
%
%============================================================================%

%---------------------------------------------------------------------------------------
%	EXPERIENCE
%----------------------------------------------------------------------------------------
\cvsection{Professional Experience}

%
\cveventthree{2016 - 2017}{Architect}{Software Defined \& Cloud Group, Hewlett Packard Enterprise (HPE)}
{Design and development of HPE's project 'New Stack' -- a continuation of the R\&D effort done under the ATG.}
{Kubernetes connector -- software networking technology that enables communication between K8s clusters.}
{Assess and set technical direction for components in 'New Stack', including API Gateway, Monitoring, and Analytics.}

%
\cveventtwo{2014 - 2016}{Principal Engineer}{Advanced Technology Group
  (ATG) for Cloud and Open Source, HPE}
{Incubate and influence new technology adoption across HPE with an
  emphasis on Cloud focused business units.}
{Kubernetes, CoreOS, nesting Container \& VM env's, System Mgmt \&
  Mon, Ansible, ELK, Load Balancing, IPv6.}

%
\cveventtwo{2011 - 2014}{Engineering Manager}{Monitoring, HP Cloud, Hewlett-Packard}
{Lead the design and operation of HP's Public Cloud (internal) monitoring and
  customer Monitoring-aaS.}
{Cloud Native \& DevOps, OpenStack( Nova, Swift, Glance, pre-Neutron
  ), Chef, collectd, RabbitMQ, Icinga, HAProxy.}

%
\cveventtwo{2005 - 2011}{Principal Engineer}{Melodeo Inc. \& HP Music, Hewlett-Packard}
{Responsible for multiple, back-end components in the Melodeo service:
  music ingestion, transcoding, and delivery.}
{Amazon-WS( EC2, S3, SQS, EMR, RDS ), Java, Python, Ruby, elastic
  transcoding, music CODEC's, Hadoop/HIVE.}

%
\cveventtwo{2003 - 2005}{Principal Engineer}{Continuous Computing}
{Developed ARM \& x86 firmware/BIOS for multiple lines of blade
  servers.  Chassis managment module dev.}
{ARM \& x86 assembly, PXE, embedded Linux \& kernel development,
  remote GDB, BusyBox, Python, GCC x-compile.}

\cveventtwo{1995 - 2003}{Software Engineer}{\it various employeers}
           {Voice over IP, CPLD Partition, Place, \& Route (CAD/EDA),
             Criminal Link Analysis, Remote Medical Coding.}
           {C++, Java, CORBA, distributed \& network programming,
             heuristic algorithms, Data Center scale networking.}

%---------------------------------------------------------------------------------------
%	Projects SECTION
%--------------------------------------------------------------------------------------
%\cvsection{Select Projects}
%
%%
%\cveventtwo{July 2014 - Present}{Lead Developer}{Glide, the Go package manager}{Developed core dependency management features which have been re-used by other projects}{Built the website, documentation, and install mechanism}
%
%%\textcolor{softcol}{\hrule}
%
%%
%\cveventone{March 2013 - Present}{Core Developer}{HTML5 Parser / Writer for PHP}{Built and shipped most popular HTML5 library for PHP}
%
%---------------------------------------------------------------------------------------
%	EDUCATION SECTION
%--------------------------------------------------------------------------------------
\cvsection{Education}

%
\cveventtwo{1989 - 1993}
           {Bachelor of Science, Computer Engineering}
           {Valparaiso University}
           {Ada language research enabling cross-host communication
             using language primitives - drove changes in ADA95.}
           {Focus on embedded system design (68k), networking, and
             VLSI design.}

\cveventtwo{1993-1995}
           {Masters of Science, Electrical Engineering {\normalfont (unconferred)} }
           {{\it Universities of Colorado \& Houston}}
           {Research:  ``Multi-threaded Processor Cache Coherency'',
             ``CPLD Design for Test'', developed Embedded Sys. course.}
           {Computer Architecture, Digital Design, Digital Signal
             Processing (DSP), CMOS/VLSI Design, Neural Networks.}

%--------------------------------------------------------------------------------------------------
%	ARTIFICIAL FOOTER (fancy footer cannot exceed linewidth)
%--------------------------------------------------------------------------------------------------

\null
\vspace*{\fill}
\hspace{-0.25\linewidth}\colorbox{bgcol}{\makebox[1.5\linewidth][c]{\mystrut
    \small \textcolor{white}{\bf www.linkedin.com/in/ericggustafson} }}


%    \textcolor{white}{www.elfwerks.org/ericg} $\cdot$
%    \textcolor{white}{github.com/egustafson}}}




%============================================================================%
%
%
%
%	DOCUMENT END
%
%
%
%============================================================================%
\end{document}
